\section{结果与讨论}

\paragraph{多类型纳米颗粒基准数据集上的模型比较。}
为获得可靠的性能基线，我们首先在大规模多类型纳米颗粒毒性数据集上比较六类模型的表现，包括随机森林\cite{breiman2001random}、XGBoost\cite{chen2016xgboost}、ABRF\cite{gromov2022attention}、LARF\cite{gromov2022larf}、TabTransformer\cite{huang2020tabtransformer} 以及蒸馏 TabTransformer\cite{hinton2015distilling}。该基准数据集来源于 Masarkar 等人建立的公开纳米颗粒毒性数据资源 \cite{masarkar2024}，覆盖金属、金属氧化物、聚合物、碳基与无机纳米颗粒，清洗后共包含 4102 个有效样本，按照 7:3 划分为训练集 3281 条与测试集 821 条。输入特征包括浓度、暴露时间、材料类别、粒径、Zeta 电位、细胞类型与检测方法等。

回归任务结果如表~\ref{tab:benchmark_regression} 所示。TabTransformer 取得最高的决定系数，$R^2$ 为 0.768，MAE 为 10.30\%，RMSE 为 14.97\%。XGBoost 与随机森林表现接近，分别达到 $R^2=0.752$ 与 $R^2=0.749$。引入注意力机制后的 ABRF 与 LARF 未超过传统树模型，蒸馏 TabTransformer 的性能则明显下降。

\begin{table}[htbp]
  \centering
  \caption{多类型纳米颗粒基准数据集上的回归性能比较}
  \label{tab:benchmark_regression}
  \begin{tabular}{lccc}
    \hline
    模型 & $R^2$ & MAE (\%) & RMSE (\%) \\
    \hline
    TabTransformer & 0.768 & 10.30 & 14.97 \\
    XGBoost & 0.752 & 11.13 & 15.56 \\
    随机森林 & 0.749 & 10.81 & 15.71 \\
    ABRF & 0.719 & 11.87 & 16.55 \\
    LARF & 0.716 & 11.94 & 16.63 \\
    蒸馏 TabTransformer & 0.506 & 16.85 & 22.31 \\
    \hline
  \end{tabular}
\end{table}

分类任务方面，树模型表现更稳健，其中 XGBoost 的 AUC 最高，为 0.945；随机森林的准确率与 F1 值最好，分别为 84.9\% 与 0.840。该基准比较表明，在大样本异质表格数据上，深度表格模型和树模型都具有较强竞争力，但树模型在分类任务中更稳定、训练成本更低。

\begin{table}[htbp]
  \centering
  \caption{多类型纳米颗粒基准数据集上的分类性能比较}
  \label{tab:benchmark_classification}
  \begin{tabular}{lcccc}
    \hline
    模型 & Accuracy & F1 & AUC \\
    \hline
    随机森林 & 0.849 & 0.840 & 0.943 \\
    XGBoost & 0.849 & 0.836 & 0.945 \\
    TabTransformer & 0.845 & 0.835 & 0.939 \\
    ABRF & 0.828 & 0.817 & 0.921 \\
    LARF & 0.819 & 0.809 & 0.914 \\
    蒸馏 TabTransformer & 0.722 & 0.707 & 0.801 \\
    \hline
  \end{tabular}
\end{table}

\paragraph{硒纳米颗粒专用数据集上的模型比较。}
在硒纳米颗粒专用数据集上，严格分析包含 434 条记录和 24 个文献组，采用 5 折外层、3 折内层的按文献分组嵌套交叉验证。基础特征包括浓度、暴露时间及其对数变换、粒径、Zeta 电位、尺寸比、表面电荷、细胞类型、修饰基团和检测方法等；化学增强特征进一步考虑元素 Se 含量、长期稳定性、结合能 Se3d 和晶型编码；IC50 设置单独作为响应相关辅助变量。所有指标均由外层 OOF 预测计算。

结果如表~\ref{tab:senp_main_results} 所示。基础随机森林的 OOF MAE 为 19.11\%、RMSE 为 24.12\%、$R^2=0.109$；XGBoost 分别为 19.49\%、23.97\% 和 0.120。不含 IC50 的化学增强没有改善严格误差：随机森林 MAE 为 19.29\%、$R^2=0.099$，XGBoost MAE 为 19.52\%、$R^2=0.128$。IC50 辅助 XGBoost 的数值指标最好（MAE 18.35\%、RMSE 23.25\%、$R^2=0.172$），但其相对基础模型的按文献配对区间跨过 0，因此只能作为响应相关辅助参考，不能解释为稳定的前瞻性收益。

\begin{table}[htbp]
  \centering
  \caption{硒纳米颗粒数据集上的严格按文献分组嵌套交叉验证结果}
  \label{tab:senp_main_results}
  \begin{tabular}{llccc}
    \hline
    模型 & 特征集 & OOF MAE (\%) & RMSE (\%) & $R^2$ \\
    \hline
    随机森林 & 基础 & 19.11 & 24.12 & 0.109 \\
    XGBoost & 基础 & 19.49 & 23.97 & 0.120 \\
    随机森林 & 化学增强，不含 IC50 & 19.29 & 24.25 & 0.099 \\
    XGBoost & 化学增强，不含 IC50 & 19.52 & 23.87 & 0.128 \\
    随机森林 & IC50 辅助 & 18.97 & 24.13 & 0.109 \\
    XGBoost & IC50 辅助 & 18.35 & 23.25 & 0.172 \\
    \hline
  \end{tabular}
\end{table}

\paragraph{剂量--时间交互消融。}
在相同的 434 条记录、24 个文献组、分折和参数预算下，单独加入 $\log(1+浓度)\times\log(1+时间)$ 交互项。相对基础配置，随机森林 MAE 增加 0.392 个百分点，按文献组 bootstrap 的 95\% 区间为 $[-0.072,1.102]$；XGBoost MAE 增加 0.605 个百分点，区间为 $[-0.199,1.016]$。两个区间均跨越 0，因此该交互项应作为阴性消融结果报告。

\paragraph{特征效应。}
当前严格仓库尚未完成针对 Modifier Category、分子量、粒径和 Zeta 电位的 SHAP 分析。Modifier Category 和分子量字段完全缺失，Zeta 电位在 434 条记录中缺失 42.40\%。因此，当前证据支持预测比较和缺失感知验证，但不足以对单个材料描述符作机制性结论。

\paragraph{关键特征对纳米硒材料选择的指导意义。}

\paragraph{组合特征分析。}
历史 31 组合搜索仍保留用于追溯，但其随机三折结果不作为严格泛化估计。当前嵌套分析在训练折内完成候选选择；不含 IC50 的化学增强没有显示稳定 MAE 降低，IC50 辅助配置仍属于响应相关结果，不能据此声称缺失的本征字段已得到验证。

\paragraph{模型结果对真实纳米硒材料设计与筛选的指导意义。}

\paragraph{回归与分类任务的最佳模型选择。}
在严格分组验证下，IC50 辅助 XGBoost 的数值回归指标最好（MAE 18.35\%、RMSE 23.25\%、$R^2=0.172$），但配对不确定性区间跨越 0。无 IC50 的随机森林基础模型（MAE 19.11\%、RMSE 24.12\%、$R^2=0.109$）仍是最清晰的部署参考，而化学增强没有改善该基线。历史基准分类结果与当前硒纳米颗粒严格回归结果应作为两条独立证据链解读。

\paragraph{与现有工作的比较。}
与 Masarkar 等人 \cite{masarkar2024} 的通用纳米颗粒毒性预测模型相比，本研究具有以下改进：
\begin{itemize}
  \item \textbf{针对单一材料（纳米硒）} 构建了专用数据集，并引入了元素 Se 含量、长期稳定性和结合能 Se3d 等材料化学增强信息，同时比较了含 IC50 与不含 IC50 的设置差异。这说明单一材料体系中的材料本征化学属性具有重要建模价值，而含 IC50 的高分数结果则更适合作为辅助参考，而非直接与仅依赖先验理化参数的通用模型进行等价比较。
  \item \textbf{系统比较了深度表格模型与树模型}，发现 TabTransformer 在小样本场景（383 条）下严重欠拟合，而树模型（尤其是 XGBoost）表现稳健。这为今后小样本纳米毒性数据集建模提供了参考。
  \item \textbf{通过穷举搜索确定了最优增强特征子集}，剔除了晶型编码等冗余特征，避免了过拟合，同时增强了模型的可解释性。
\end{itemize}
与近期其他基于注意力机制的树模型（如 ABRF、LARF）相比，本研究未发现注意力机制在小规模表格数据上带来显著提升，反而增加了计算开销。因此，对于纳米毒性这类中等规模、混合型特征的数据，传统的梯度提升树（XGBoost）仍是性价比最高的选择。

\paragraph{结论。}
严格重建表明，按文献分组后模型表现明显低于历史随机划分结果。在当前 434 条归档数据上，化学增强和显式剂量--时间交互均未产生稳定改善；IC50 辅助 XGBoost 只能作为响应相关敏感性参考，不能作为前瞻模型。后续外部验证必须在新 SeNP 湿实验前冻结预测，补齐材料表征，并使用不把实验 IC50 作为输入的独立剂量--反应数据。
